\begin{examplebox}
	\textbf{\textcolor{CExample}{例 1（基础）}}

	\textbf{\textcolor{CExample}{题目：}}
	已知正四面体棱长为$4$，求高$h$和内切球半径$r$。

	\textbf{\textcolor{CExample}{解题步骤：}}
	\begin{description}[style=nextline, leftmargin=2.8em, labelsep=0.5em, itemsep=0.45em]
		\item[\textbf{第一步（套用高公式）：}] $h = \frac{\sqrt{6}}{3} \times 4 = \frac{4\sqrt{6}}{3}$。
			\par\textbf{理由：}正四面体高公式$h = \frac{\sqrt{6}}{3}a$，代入$a=4$。
		\item[\textbf{第二步（套用球半径公式）：}] $r = \frac{\sqrt{6}}{12} \times 4 = \frac{\sqrt{6}}{3}$。
			\par\textbf{理由：}内切球半径公式$r = \frac{\sqrt{6}}{12}a$，代入$a=4$。
		\item[\textbf{第三步（比例验证）：}] $r = \frac{h}{4} = \frac{1}{4}\cdot\frac{4\sqrt{6}}{3} = \frac{\sqrt{6}}{3}$，吻合。
			\par\textbf{理由：}验证$r = h/4$，确保计算正确。
	\end{description}

	\textbf{\textcolor{CExample}{关键结论：}}
	\textbf{$h = \frac{4\sqrt{6}}{3}$，$r = \frac{\sqrt{6}}{3}$。}

	\medskip
	\textbf{\textcolor{CExample}{例 2（稍变形）}}

	\textbf{\textcolor{CExample}{题目：}}
	已知正四面体内切球半径为$2$，求棱长$a$和高$h$。

	\textbf{\textcolor{CExample}{解题步骤：}}
	\begin{description}[style=nextline, leftmargin=2.8em, labelsep=0.5em, itemsep=0.45em]
		\item[\textbf{第一步（反求棱长）：}] 由$r = \frac{\sqrt{6}}{12}a$得$a = \frac{12r}{\sqrt{6}} = \frac{12\times 2}{\sqrt{6}} = \frac{24}{\sqrt{6}} = 4\sqrt{6}$。
			\par\textbf{理由：}解出$a$。
		\item[\textbf{第二步（求高）：}] $h = \frac{\sqrt{6}}{3}a = \frac{\sqrt{6}}{3} \times 4\sqrt{6} = \frac{6\times 4}{3} = 8$。
			\par\textbf{理由：}代入$a$计算。
		\item[\textbf{第三步（合理性验证）：}] $r = \frac{h}{4} = \frac{8}{4} = 2$，与已知一致。
			\par\textbf{理由：}比例验证。
	\end{description}

	\textbf{\textcolor{CExample}{关键结论：}}
	\textbf{$a = 4\sqrt{6}$，$h = 8$。}

	\medskip
	\textbf{\textcolor{CExample}{例 3（综合应用）}}

	\textbf{\textcolor{CExample}{题目：}}
	已知正四面体的高为$6$，求其内切球半径$r$、球心到底面距离$d$以及棱长$a$。

	\textbf{\textcolor{CExample}{解题步骤：}}
	\begin{description}[style=nextline, leftmargin=2.8em, labelsep=0.5em, itemsep=0.45em]
		\item[\textbf{第一步（求内切球半径）：}] $r = \frac{h}{4} = \frac{6}{4} = 1.5$。
			\par\textbf{理由：} $r = h/4$。
		\item[\textbf{第二步（球心到底面距离）：}] $d = r = 1.5$。
			\par\textbf{理由：}内切球球心到各面距离均为$r$，故到底面距离等于$r$。
		\item[\textbf{第三步（反求棱长）：}] 由$h = \frac{\sqrt{6}}{3}a$得$a = \frac{3h}{\sqrt{6}} = \frac{18}{\sqrt{6}} = 3\sqrt{6}$。
			\par\textbf{理由：}解出$a$。
	\end{description}

	\textbf{\textcolor{CExample}{关键结论：}}
	\textbf{$r = 1.5$，$d = 1.5$，$a = 3\sqrt{6}$。}
\end{examplebox}
